% ============================================================================
% Fichier  : annexes/Ann_E_bibliographie.tex
% Titre    : Bibliographie Commentée
% Version  : 1.0 -- novembre 2026
% ============================================================================
% Les textes de l'auteur sont des prépublications ; les travaux en
% préparation (A2 à A5) ne sont pas cités (cadre de cohérence du corpus).
% ============================================================================

\chapter{Bibliographie Commentée}
\label{ann:bibliographie}

Pour chaque partie, quelques lectures, avec ce qu'elles apportent. Les
références complètes figurent dans la bibliographie du livre. Les textes de
l'auteur sont des prépublications~; les travaux en préparation ne sont pas
cités.

\section*{Pour commencer}

\begin{itemize}
    \item \textbf{Le programme de l'auditologie}~\cite{minka2026auditologie}~:
          objet, axiomes, piliers et méthode~; à lire après le prologue, et à
          relire avec le Chapitre~\ref{chap:auditologie}.
    \item \textbf{L'instrument trichotomique}~\cite{minka2026nek}~: la
          réponse \NEK{}, ses conditions, les mesures et l'indice d'usage
          stratégique.
\end{itemize}

\section*{Partie I · Fondements}

\begin{itemize}
    \item L'IIA, le modèle des trois lignes~\cite{iia2020troislignes}~: la
          répartition des rôles, en une dizaine de pages.
    \item L'ISO/IEC~27014~\cite{iso27014_2020} et la fonction \emph{Gouverner}
          du NIST CSF~2.0~\cite{nist2024csf}~: la gouvernance de la sécurité.
    \item L'ISO~31000~\cite{iso31000_2018}~: le vocabulaire du risque.
    \item Cox, sur les matrices de risque~\cite{cox2008matrices}~: ce que
          les couleurs cachent.
\end{itemize}

\section*{Partie II · Méthodes}

\begin{itemize}
    \item L'ISO/IEC~27005:2022~\cite{iso27005_2022}~: le processus appliqué à
          la sécurité de l'information.
    \item Le guide EBIOS Risk Manager de l'ANSSI~\cite{anssi2018ebios} et ses
          fiches méthodes~: la méthode et ses ateliers.
    \item Freund et Jones~\cite{freund2015fair}~: l'exposé de référence de
          FAIR~; Hubbard et Seiersen~\cite{hubbard2016measure}~: le calibrage
          des experts.
    \item L'ISO~19011~\cite{iso19011_2018} et le NIST SP~800-53A~\cite{nist2022sp80053a}~:
          les principes de l'audit et les méthodes d'évaluation.
\end{itemize}

\section*{Partie III · Cadre normatif et juridique}

\begin{itemize}
    \item L'ISO/IEC~27001:2022~\cite{iso27001_2022} et l'ISO/IEC~27002:2022~\cite{iso27002_2022}~:
          à lire ensemble, l'une pour les exigences, l'autre pour les
          mesures.
    \item Les CIS Controls~\cite{cis2021controls}~: une hygiène de base
          priorisée.
    \item La loi n°~2010/012~\cite{loi2010012}, dans son texte officiel, et le
          règlement COBAC R-2016/04~\cite{cobac2016r04}.
    \item Les normes de l'IIA~\cite{iia2024gias} et le code de
          l'IESBA~\cite{iesba2023code}~: la déontologie.
\end{itemize}

\section*{Partie IV · Pratique}

\begin{itemize}
    \item L'ISO/IEC~27007~\cite{iso27007_2020} et l'ISO/IEC~27004~\cite{iso27004_2016}~:
          auditer et mesurer un SMSI.
    \item Les CIS Benchmarks~\cite{cis2026benchmarks} et la spécification
          CVSS~4.0~\cite{first2023cvss4}~: l'audit technique.
\end{itemize}

\section*{Parties V et VI · Défis avancés et recherche}

\begin{itemize}
    \item L'ISO/IEC~42001~\cite{iso42001_2023} et le NIST AI RMF~\cite{nist2023airmf}~:
          les cadres de l'IA.
    \item L'ISO/IEC~27017~\cite{iso27017_2015}~: la sécurité du nuage.
    \item Shannon~\cite{shannon1948}, Kullback et Leibler~\cite{kullback1951},
          Lakatos~\cite{lakatos1970}~: les fondements des mesures et de la
          méthode.
\end{itemize}
